
\documentclass{article}
\usepackage{colortbl}
\usepackage{makecell}
\usepackage{multirow}
\usepackage{supertabular}

\begin{document}

\newcounter{utterance}

\centering \large Interaction Transcript for game `codenames', experiment `opponent\_difficult', episode 3 with reflexion{-}Qwen3.5{-}27B.
\vspace{24pt}

{ \footnotesize  \setcounter{utterance}{1}
\setlength{\tabcolsep}{0pt}
\begin{supertabular}{c@{$\;$}|p{.15\linewidth}@{}p{.15\linewidth}p{.15\linewidth}p{.15\linewidth}p{.15\linewidth}p{.15\linewidth}}
    \# & $\;$A & \multicolumn{4}{c}{Game Master} & $\;\:$B\\
    \hline

    \theutterance \stepcounter{utterance}  
    & & & \multicolumn{2}{p{0.3\linewidth}}{
        \cellcolor[rgb]{0.9,0.9,0.9}{
            \makecell[{{p{\linewidth}}}]{
                \texttt{\tiny{[GM$|$GM]}}
                {'hidden': {'team': ['ring', 'wall'], 'innocent': ['field', 'wind', 'europe', 'snow', 'shop'], 'opponent': [], 'assassin': ['pin']}, 'revealed': {'team': {'team': [], 'innocent': [], 'opponent': [], 'assassin': []}, 'opponent': {'team': [], 'innocent': [], 'opponent': [], 'assassin': []}}}
            }
        }
    }
    & & \\ \\

    \theutterance \stepcounter{utterance}  
    & & \multicolumn{4}{p{0.6\linewidth}}{
        \cellcolor[rgb]{0.9,0.9,0.9}{
            \makecell[{{p{\linewidth}}}]{
                \texttt{\tiny{[P1$\langle$GM]}}
                \texttt{You have attempted this same task before and did not succeed. These are the plans you wrote after those attempts.} \\
\texttt{Trial \#1: New plan: My previous strategy failed because I provided clues that were too broad or ambiguous, leading my partner to guess distractor words (like "lock" for "steal" and "watch" for "view") or fail to guess the intended targets (like "wall" for "boxing"). Specifically, "boxing" is a weak link for "wall" (which is likely a "brick wall" or "boundary" in this context, not a boxing ring wall), and "steal" directly triggered the distractor "lock". Next time, I will prioritize clues that create a unique, strong semantic link to the target while explicitly avoiding words that share strong associations with the remaining distractors. For "ring" and "wall", I should have used a clue like "jewelry" (for ring) or "boundary" (for wall) to avoid the "boxing" ambiguity, or split them into separate turns to ensure accuracy. I will also verify that my clue does not trigger any remaining distractor words before submitting.} \\
\\ 
\texttt{Trial \#2: New plan: My previous attempt failed because I only targeted "ring" with the clue "jewelry," leaving "wall" unguessed and failing to maximize the efficiency of the turn. Since the goal is to guess all team words, I must provide a clue that connects to both "ring" and "wall" simultaneously if possible, or ensure I address all remaining words in the fewest turns. A better strategy is to use a clue like "boundary" or "enclosure," which semantically links to both a "ring" (as in a boxing ring or arena) and a "wall" (as a barrier), allowing the partner to guess both words in one turn. If a single clue for both is too risky, I will prioritize a clue that covers the most words without triggering distractors like "field" or "wind," rather than isolating a single word and leaving others for a subsequent turn that might not happen.} \\
\\ 
\\ 
\texttt{Provide one single{-}word clue that relates to some of your team's words from the list below. You can choose to how many and to which words your clue relates to.} \\
\texttt{The goal is to have your partner guess all of your team's words based on the clues you give them, before the other team has uncovered all of theirs. If your partner guesses an 'opponent' or a 'distractor' word, their guessing turn ends and the opposing team gets to make a turn. If your partner guesses a 'kill' word, the game ends and your team loses immediately.} \\
\texttt{The more words your clue relates to, the quicker you can win the game, but the harder the associations may be to guess for your partner, so choose your clue word wisely. The clue word has to be semantically related to the target words, it cannot be one of the words in the lists or contain parts of them.} \\
\\ 
\texttt{Always give your single{-}word clue and your comma{-}separated list of related target words in the following format and make your answers as short as possible, never include any other text than is required in this form:} \\
\\ 
\texttt{CLUE: <WORD>} \\
\texttt{TARGETS: <TARGETS>} \\
\\ 
\texttt{Your team words are: ring, wall.} \\
\texttt{Your opponent's team words are: .} \\
\texttt{Distractor words are: field, wind, europe, snow, shop.} \\
\texttt{The kill words are: pin.} \\
            }
        }
    }
    & \\ \\

    \theutterance \stepcounter{utterance}  
    & \multicolumn{4}{p{0.6\linewidth}}{
        \cellcolor[rgb]{0.9,0.9,0.9}{
            \makecell[{{p{\linewidth}}}]{
                \texttt{\tiny{[P1$\rangle$GM]}}
                \texttt{CLUE: boundary} \\
\texttt{TARGETS: ring, wall} \\
            }
        }
    }
    & & \\ \\

    \theutterance \stepcounter{utterance}  
    & & & \multicolumn{2}{p{0.3\linewidth}}{
        \cellcolor[rgb]{0.9,0.9,0.9}{
            \makecell[{{p{\linewidth}}}]{
                \texttt{\tiny{[GM$|$GM]}}
                \texttt{boundary} \\
            }
        }
    }
    & & \\ \\

    \theutterance \stepcounter{utterance}  
    & & & \multicolumn{2}{p{0.3\linewidth}}{
        \cellcolor[rgb]{0.9,0.9,0.9}{
            \makecell[{{p{\linewidth}}}]{
                \texttt{\tiny{[GM$|$GM]}}
                ['ring', 'wall']
            }
        }
    }
    & & \\ \\

    \theutterance \stepcounter{utterance}  
    & & & \multicolumn{2}{p{0.3\linewidth}}{
        \cellcolor[rgb]{0.9,0.9,0.9}{
            \makecell[{{p{\linewidth}}}]{
                \texttt{\tiny{[GM$|$GM]}}
                {'word': 'ring', 'assignment': 'team'}
            }
        }
    }
    & & \\ \\

    \theutterance \stepcounter{utterance}  
    & & & \multicolumn{2}{p{0.3\linewidth}}{
        \cellcolor[rgb]{0.9,0.9,0.9}{
            \makecell[{{p{\linewidth}}}]{
                \texttt{\tiny{[GM$|$GM]}}
                {'word': 'wall', 'assignment': 'team'}
            }
        }
    }
    & & \\ \\

    \theutterance \stepcounter{utterance}  
    & & & \multicolumn{2}{p{0.3\linewidth}}{
        \cellcolor[rgb]{0.9,0.9,0.9}{
            \makecell[{{p{\linewidth}}}]{
                \texttt{\tiny{[GM$|$GM]}}
                \texttt{opponent has won} \\
            }
        }
    }
    & & \\ \\

    \theutterance \stepcounter{utterance}  
    & \multicolumn{4}{p{0.6\linewidth}}{
        \cellcolor[rgb]{0.95,0.95,0.95}{
            \makecell[{{p{\linewidth}}}]{
                \texttt{\tiny{[P2$\rightarrow$P1]}}
                \texttt{Post{-}episode Reflexion (not sent to GM):} \\
\texttt{New plan: My previous attempt failed because the clue "boundary" was too abstract, causing the partner to likely guess a distractor like "field" (which has boundaries) or fail to connect the specific concepts of "ring" and "wall" quickly enough. The strategy of grouping both words under a broad concept backfired. Next time, I will prioritize a clue that creates a concrete, shared physical context for both "ring" and "wall" that is distinct from the distractors. I will use the clue "fence," as it strongly implies a "ring" (enclosure) and a "wall" (barrier) while being semantically distant from "field," "wind," "europe," "snow," "shop," and the kill word "pin." If a single strong link is not immediately obvious, I will split the targets, guessing one word per turn with a highly specific clue (e.g., "jewelry" for "ring") to ensure at least one point is secured without risking a turn{-}ending distractor guess.} \\
            }
        }
    }
    & & \\ \\

\end{supertabular}
}

\end{document}
